\documentclass[11pt,a4paper]{article}
\usepackage[T1]{fontenc}
\usepackage{lmodern}
\usepackage[margin=22mm]{geometry}
\usepackage{booktabs,tabularx,longtable,enumitem,xcolor,hyperref}
\hypersetup{colorlinks=true,urlcolor=blue!50!black}
\setlength{\parindent}{0pt}\setlength{\parskip}{6pt}
\setlist{nosep,leftmargin=*}
\newcommand{\entry}[1]{\textcolor{blue!60!black}{\textit{[#1]}}}
\newcommand{\code}[1]{\texttt{\detokenize{#1}}}
\title{DAPHNE--SSP scan report\\\large Editable operator template}
\author{\entry{Operator and institute}}
\date{\entry{Date, time range, and time zone}}
\begin{document}\maketitle
\textbf{Instructions.} Replace bracketed fields with observations and artifact paths. Delete instructions before circulation. Mark a check as failed or not tested when evidence is absent. The reference run 45371 belongs to a separate completed scan; do not reuse its numbers as a new measurement.

\section{Purpose and result}
\entry{One paragraph: what was scanned, which channels were acquired, whether it completed, and the principal finding. Distinguish an electronics/control exercise from a detector gain measurement.}

\begin{tabularx}{\linewidth}{@{}lX@{}}\toprule
Identity & Recorded value \\\midrule
Board / location & \entry{Asset ID; rack; installation or cabling changes} \\
Operator / control host & \entry{Authorized operating account and host} \\
DAQ run / scan tag & \entry{Actual DAQ run number; unique scan ID} \\
Campaign directory & \entry{Absolute persistent path} \\
Gateware profile / identity & \entry{Active profile; live build ID; source commit} \\
Software / environment & \entry{Server; PDS; decoder; DAQ release; local changes} \\
Evidence manifest & \entry{Path to SHA-256 manifest and source/dirty-state record} \\\bottomrule
\end{tabularx}

\section{Physical and logical configuration}
\entry{List management MAC/IP and actual data MAC/IP values separately from registered reservations. State any address mismatch and the DAQ configuration used.}

\begin{tabularx}{\linewidth}{@{}lX@{}}\toprule
Setting & Recorded value \\\midrule
Links / senders / channels & \entry{Physical links; wire stream IDs; DAQ source IDs; selected channels} \\
AFE mapping & \entry{Firmware AFE to physical board AFE mapping} \\
DAPHNE endpoint / master & \entry{Address; master name; partition if relevant; state; locks; timestamp validity} \\
SSP endpoint / master & \entry{Control address; timing address; master; state; readiness} \\
Source configuration & \entry{For each master: generator, command ID, requested and measured rate, period, clock} \\
DAQ window / trigger rate & \entry{Readout window in seconds and DAQ request rate in Hz} \\
Optical/timing relationship & \entry{Optically coupled? Shared master? How verified?} \\\bottomrule
\end{tabularx}

\section{Scan settings and method}
\begin{tabularx}{\linewidth}{@{}lX@{}}\toprule
Control & Recorded value and native units \\\midrule
Offset / VGAIN & \entry{DAC offset; VGAIN points; common to all AFEs or sequential?} \\
Detector bias & \entry{DAPHNE detector-bias setting; separate from SSP LED bias} \\
SSP LED selection & \entry{Controller channel count; mask; enabled channels} \\
SSP LED bias & \entry{270 nm and 367 nm native codes} \\
SSP pulse settings & \entry{Single/double mode; widths in ticks; delay in ticks; burst count} \\
Per-point selection & \entry{Window duration; target count if used; settling; guard/tail} \\
Final-state request & \entry{Paused or explicitly configured operating state} \\\bottomrule
\end{tabularx}

\entry{Describe the per-point sequence: pause/drain, DAC writes, SSP configuration/readback, timing starts, acquisition, counter snapshots, stops, and cleanup. Include exact launcher command and script/configuration hashes.}

\section{Results and acceptance}
\textbf{Denominators.} Selected analysis windows contain \entry{waveforms per channel per point and total samples}. Hardware counters cover \entry{time interval and accepted master commands per channel}. These are different observation intervals; report both.

\begin{center}\small
\begin{tabular}{@{}rrrrr@{}}\toprule
Point / code & Frames/channel & Pedestal range & Within-frame RMS & Rail clipping \\
 & & ADC counts & ADC counts & zero / full scale \\\midrule
\entry{point} & \entry{count} & \entry{min--max} & \entry{min--max} & \entry{fractions} \\
\entry{point} & \entry{count} & \entry{min--max} & \entry{min--max} & \entry{fractions} \\
\entry{point} & \entry{count} & \entry{min--max} & \entry{min--max} & \entry{fractions} \\
\bottomrule\end{tabular}
\end{center}

\entry{Attach pedestal/RMS versus scan coordinate and representative waveform figures. Define pedestal, total RMS, within-frame RMS, and waveform-mean drift. State any rejection, baseline subtraction, or filtering.}

\begin{longtable}{@{}p{0.51\linewidth}p{0.43\linewidth}@{}}\toprule
Check & Pass/fail/not tested, and evidence \\\midrule\endhead
Expected gateware identity and runtime health & \entry{Observed identity and status log} \\
Endpoint clock lock and timestamp readiness & \entry{Before/after snapshots for both devices} \\
Every selected source/channel present & \entry{Source IDs, channel counts, receipt paths} \\
Common DAPHNE event timestamps; no gaps/duplicates & \entry{Windows and audit result} \\
Fragment status / format / stream mapping & \entry{Status bits and decoder checks} \\
Hardware records versus accepted master commands & \entry{Per-channel increments and master denominator} \\
Busy/full drops and rejected commands & \entry{Counts and observation interval} \\
SSP bias/duration readbacks & \entry{Requested versus read-back values} \\
SSP command counter response & \entry{Counter delta; SSP master; no cross-domain inference} \\
Input hashes and official decoder cross-check & \entry{Manifest and decoder revision} \\
ADC rail clipping and anomalous channels & \entry{Fractions; affected channels/points} \\
\bottomrule\end{longtable}

\section{Final state, limitations, and next action}
\entry{Record whether DAQ stop-run succeeded, whether applications are in error, DAPHNE pause mask, stopped timing sources, SSP LED-disable readback, and final endpoint readiness. Include cleanup failures explicitly.}

\entry{State what was not established: optical output, detector response, physical gain, shared-master synchronization, reboot persistence, or other configuration limits. Describe the next test, owner, and required setup.}

\section{Artifacts and review}
\begin{itemize}
\item Raw HDF5 files, record IDs, windows, and hashes: \entry{paths/manifest}.
\item Selected waveform binaries and receipts: \entry{paths/hashes}.
\item Plan, runner, acquisition adapter, SSP driver and routine: \entry{versions/hashes/dirty state}.
\item Numeric analysis, plots, and decoder provenance: \entry{paths/hashes}.
\item Failed/interrupted attempts retained separately: \entry{IDs and disposition}.
\item Reviewer / date / status: \entry{not reviewed, reviewed, or approved; name the authority if approved}.
\end{itemize}
\end{document}
